% ============================================================
%  Pin Mapping and Configuration Reference
%  Raspberry Pi Zero 2 W  --- Colorimeter Prototype
%  As-Built Document  ---  Pre-Testing Reference
%
%  Components covered:
%    Adafruit 1651  2.8" TFT Touch Shield (ILI9341 + TSC2007)
%    74HC595 8-bit shift register (SPI1)
%    AO3400A N-channel MOSFETs x10 (8 via 595, 2 direct GPIO)
%    CEM-1203(42) passive buzzer via MMBT3904 NPN BJT
%    1-Wire temperature probe (non-default GPIO16)
%    Connectors J1--J8, jumpers JP1--JP2
%
%  Sources:
%    RPi GPIO: raspberrypi.com/documentation/computers/raspberry-pi.html#gpio
%    Adafruit 1651: learn.adafruit.com/adafruit-2-8-tft-touch-shield-v2
%    74HC595: TI SN74HC595 datasheet; Adafruit learn.adafruit.com/74hc595
%    AO3400A: Alpha & Omega AO3400A datasheet
%    pigpio: abyz.me.uk/rpi/pigpio/
% ============================================================
\documentclass[11pt,a4paper]{article}

\usepackage[a4paper, margin=1.8cm]{geometry}
\usepackage{booktabs}
\usepackage{tabularx}
\usepackage{longtable}
\usepackage{array}
\usepackage{xcolor}
\usepackage{colortbl}
\usepackage{hyperref}
\usepackage{parskip}
\usepackage{microtype}
\usepackage{fancyhdr}
\usepackage{titlesec}
\usepackage{mdframed}
\usepackage{enumitem}
\usepackage{multirow}
\usepackage{amsmath}
\usepackage{alltt}

% ── Colours ──────────────────────────────────────────────────
\definecolor{hdr}{RGB}{30,70,120}
\definecolor{rowB}{RGB}{248,251,255}
\definecolor{warn}{RGB}{255,243,205}
\definecolor{warnborder}{RGB}{200,150,0}
\definecolor{notebg}{RGB}{237,247,237}
\definecolor{noteborder}{RGB}{60,140,60}
\definecolor{oddcol}{RGB}{220,235,255}
\definecolor{evencol}{RGB}{220,250,225}
\definecolor{pwrcol}{RGB}{255,230,220}
\definecolor{errcol}{RGB}{255,210,210}

% ── Column types ─────────────────────────────────────────────
\newcolumntype{L}[1]{>{\raggedright\arraybackslash}p{#1}}
\newcolumntype{C}[1]{>{\centering\arraybackslash}p{#1}}

% ── Section styling ───────────────────────────────────────────
\titleformat{\section}{\large\bfseries\color{hdr}}{}{0em}{}[\titlerule]
\titleformat{\subsection}{\normalsize\bfseries\color{hdr}}{}{0em}{}
\titleformat{\subsubsection}{\normalsize\itshape\color{hdr}}{}{0em}{}

% ── Header / footer ───────────────────────────────────────────
\pagestyle{fancy}
\fancyhf{}
\lhead{\small Colorimeter Prototype --- RPi Zero 2W Pin Mapping and Configuration}
\rhead{\small As-Built Pre-Testing Reference}
\cfoot{\small Page \thepage}
\renewcommand{\headrulewidth}{0.4pt}

% ── Coloured boxes ────────────────────────────────────────────
\newmdenv[backgroundcolor=notebg, linecolor=noteborder, linewidth=1.2pt,
  leftmargin=0pt, rightmargin=0pt, innerleftmargin=8pt, innerrightmargin=8pt,
  innertopmargin=6pt, innerbottommargin=6pt,
  skipabove=6pt, skipbelow=6pt]{notebox}

\newmdenv[backgroundcolor=warn, linecolor=warnborder, linewidth=1.2pt,
  leftmargin=0pt, rightmargin=0pt, innerleftmargin=8pt, innerrightmargin=8pt,
  innertopmargin=6pt, innerbottommargin=6pt,
  skipabove=6pt, skipbelow=6pt]{warnbox}

\newmdenv[backgroundcolor=errcol, linecolor=red, linewidth=1.5pt,
  leftmargin=0pt, rightmargin=0pt, innerleftmargin=8pt, innerrightmargin=8pt,
  innertopmargin=6pt, innerbottommargin=6pt,
  skipabove=6pt, skipbelow=6pt]{errorbox}

% ── Helpers ───────────────────────────────────────────────────
\newcommand{\fixed}{\textbf{FIXED}}
\newcommand{\soft}{\textit{soft}}
\newcommand{\na}{---}
\newcommand{\jref}[1]{\texttt{\small[#1]}}

% =============================================================
\begin{document}

\begin{center}
  {\LARGE\bfseries\color{hdr} Colorimeter Prototype\\[4pt]
   Pin Mapping and Configuration Reference}\\[6pt]
  {\large Raspberry Pi Zero 2~W --- As-Built}\\[4pt]
  {\normalsize Pre-testing reference document.
   Describes all hardware connections, required OS configuration,\\
   known issues, and testing cautions as of first prototype build.}\\[4pt]
  {\small\color{gray}
   RPi GPIO reference:
   \href{https://www.raspberrypi.com/documentation/computers/raspberry-pi.html\#gpio}%
   {raspberrypi.com/documentation/.../raspberry-pi.html\#gpio}\\
   Adafruit 1651:
   \href{https://learn.adafruit.com/adafruit-2-8-tft-touch-shield-v2}%
   {learn.adafruit.com/adafruit-2-8-tft-touch-shield-v2}\\
   74HC595: TI SN74HC595 datasheet;
   \href{https://learn.adafruit.com/74hc595}{learn.adafruit.com/74hc595}\\
   pigpio library:
   \href{http://abyz.me.uk/rpi/pigpio/}{abyz.me.uk/rpi/pigpio}}
\end{center}
\vspace{4pt}\hrule\vspace{8pt}

% =============================================================
\section{RPi Zero 2~W GPIO Peripheral Reference}
\label{sec:gpioref}

All GPIO pins operate at \textbf{3.3\,V logic}.
Applying 5\,V to any GPIO pin will permanently damage the SoC.
The following peripheral pin assignments are \textbf{fixed by the RP3A0 SoC}
and cannot be remapped in software.

\begin{notebox}
\textbf{Source:}
\href{https://www.raspberrypi.com/documentation/computers/raspberry-pi.html\#gpio}%
{raspberrypi.com --- GPIO and 40-pin header documentation}

\medskip
\textbf{Hardware PWM} --- available on four pins only:
GPIO12 (PWM0-A, pin~32), GPIO13 (PWM1, pin~33),
GPIO18 (PWM0-B, pin~12), GPIO19 (PWM0, pin~35).
\textit{GPIO12 and GPIO13 share one clock divider.
GPIO18 and GPIO19 share another.
Both channels within a pair can have independent duty cycles
but their frequencies must be derived from the same base clock.}

\medskip
\textbf{SPI0} (enabled by \texttt{dtparam=spi=on}):
MOSI=GPIO10/pin~19; MISO=GPIO9/pin~21; SCLK=GPIO11/pin~23;
CE0=GPIO8/pin~24; CE1=GPIO7/pin~26.

\medskip
\textbf{SPI1} (enabled by \texttt{dtoverlay=spi1-1cs}):
MOSI=GPIO20/pin~38; MISO=GPIO19/pin~35; SCLK=GPIO21/pin~40;
CE0=GPIO18/pin~12; CE1=GPIO17/pin~11; CE2=GPIO16/pin~36.
\textit{Warning: the spi1 overlay claims GPIO18 as CE0 by default.
GPIO18 is used in this design for a separate purpose.
See Section~\ref{sec:spi1config} for correct configuration.}

\medskip
\textbf{I2C bus~1} (enabled by \texttt{dtparam=i2c\_arm=on}):
SDA=GPIO2/pin~3; SCL=GPIO3/pin~5.

\medskip
\textbf{UART}:
TX=GPIO14/pin~8; RX=GPIO15/pin~10.
\textit{Reserved for Linux serial console. Do not use.}

\medskip
\textbf{EEPROM}: GPIO0/pin~27 and GPIO1/pin~28.
\textit{Physical pins removed on this prototype. Do not use.}
\end{notebox}

% ── Power notes ───────────────────────────────────────────────
\subsection{Power and Ground Pin Notes}

\begin{warnbox}
\textbf{5\,V pins --- Pin~2 and Pin~4 (even side):}
Electrically identical. Both connect directly to the micro-USB supply rail.
In this design: Pin~2 supplies the Adafruit 1651 shield via J1.
Pin~4 supplies the LED load circuit via J5--J8.
Current limit is set by the micro-USB supply and the board polyfuse
($\approx$2\,A for the Zero 2\,W).

\smallskip
\textbf{3.3\,V pins --- Pin~1 (odd) and Pin~17 (odd):}
Both are outputs of the same onboard 3.3\,V regulator. Electrically identical.
\textbf{Total current budget across the entire 3.3\,V rail is $\approx$50\,mA},
shared among GPIO output drivers and all 3.3\,V peripherals.
In this design the 3.3\,V rail supplies: 74HC595 VCC, 1-Wire pull-up,
I2C pull-ups (via 1651 shield), and J2 pin~1 (1-Wire device supply).
The 1651 shield has its own onboard 3.3\,V regulator fed from 5\,V ---
it does \emph{not} draw from the RPi 3.3\,V rail.

\smallskip
\textbf{Ground:} Pins 6, 9, 14, 20, 25, 30, 34, 39 are all identical ground.
Use the ground pin nearest to each signal group for lowest-impedance return paths.
\end{warnbox}

% =============================================================
\section{Complete RPi Zero 2~W Pin Allocation --- As-Built}
\label{sec:pinmap}

\textbf{Legend:}
\colorbox{oddcol}{\strut Odd-side pin (1651 / SPI0 / I2C)} \quad
\colorbox{evencol}{\strut Even-side pin (general I/O)} \quad
\colorbox{pwrcol}{\strut Power / Ground} \quad
\colorbox{errcol}{\strut Note / caution}

\smallskip
\textbf{Columns:} Phys = physical header pin;
BCM = Broadcom GPIO number (use in Python);
Dir = O output / I input / IO bidirectional / PWM / PWR;
Connector reference in \texttt{[brackets]}.

\vspace{4pt}
\setlength{\tabcolsep}{3pt}
\renewcommand{\arraystretch}{1.35}

\begin{longtable}{C{0.85cm} C{1.2cm} C{0.65cm} L{3.2cm} L{8.1cm}}
  \toprule
  \rowcolor{hdr}
  \color{white}\textbf{Phys} &
  \color{white}\textbf{BCM} &
  \color{white}\textbf{Dir} &
  \color{white}\textbf{Assignment} &
  \color{white}\textbf{Notes / connector} \\
  \midrule
  \endfirsthead
  \toprule
  \rowcolor{hdr}
  \color{white}\textbf{Phys} &
  \color{white}\textbf{BCM} &
  \color{white}\textbf{Dir} &
  \color{white}\textbf{Assignment} &
  \color{white}\textbf{Notes / connector} \\
  \midrule
  \endhead
  \bottomrule
  \endfoot

  % ── ODD SIDE ────────────────────────────────────────────────
  \multicolumn{5}{l}{\cellcolor{oddcol}%
    \textbf{Odd-side pins --- Adafruit 1651 display/touch shield
    (SPI0, I2C, PWM backlight)}} \\
  \midrule

  \rowcolor{pwrcol}
  1 & \na & PWR & 3.3\,V &
    Not connected to 1651 (shield has own regulator).
    Available for low-current 3.3\,V loads ($\leq$50\,mA total rail). \\

  \rowcolor{oddcol}
  3 & GPIO2 & IO & I2C SDA $\to$ TSC2007 &
    \fixed\ by SoC. I2C bus~1 data.
    Shield has onboard 4.7\,k$\Omega$ pull-up --- do not add PCB pull-up.
    33\,$\Omega$ series ESD resistor on PCB.
    \jref{J1-Pin2} \\

  \rowcolor{rowB}
  5 & GPIO3 & IO & I2C SCL $\to$ TSC2007 &
    \fixed\ by SoC. I2C bus~1 clock.
    Same pull-up note as SDA.
    33\,$\Omega$ series ESD resistor on PCB.
    \jref{J1-Pin4} \\

  \rowcolor{oddcol}
  7 & GPIO4 & I & TSIRQ --- touch interrupt &
    Active-low interrupt output from TSC2007.
    Enable RPi internal pull-up in software:
    \texttt{GPIO.setup(4, GPIO.IN, pull\_up\_down=GPIO.PUD\_UP)}.
    33\,$\Omega$ series ESD resistor on PCB.
    \jref{J1-Pin10} \\

  \rowcolor{pwrcol}
  9 & \na & GND & Ground &
    Odd-side ground return. Connects to J1-Pin3, J1-Pin5, J1-Pin15.
    \jref{J1-Pin3, J1-Pin5, J1-Pin15} \\

  \rowcolor{oddcol}
  11 & GPIO17 & \na & Spare &
    Unassigned. Routed to J1 for future use.
    33\,$\Omega$ series resistor on PCB.
    \jref{J1-Pin14} \\

  \rowcolor{pwrcol}
  13 & GPIO27 & \na & No connect &
    Unassigned. Physical pin present but no PCB connection. \\

  \rowcolor{pwrcol}
  15 & GPIO22 & \na & No connect &
    Unassigned. Physical pin present but no PCB connection. \\

  \rowcolor{pwrcol}
  17 & \na & PWR & 3.3\,V &
    Not connected to 1651. Same rail as pin~1.
    Available for low-current 3.3\,V loads. \\

  \rowcolor{oddcol}
  19 & GPIO10 & O & SPI0 MOSI $\to$ ILI9341 &
    \fixed\ by SoC (SPI0 MOSI). High-speed serial data to display controller.
    33\,$\Omega$ series resistor on PCB.
    \jref{J1-Pin13} \\

  \rowcolor{rowB}
  21 & GPIO9 & I & SPI0 MISO $\leftarrow$ ILI9341 &
    \fixed\ by SoC (SPI0 MISO). Read-back from display and optional SD card.
    33\,$\Omega$ series resistor on PCB.
    \jref{J1-Pin12} \\

  \rowcolor{oddcol}
  23 & GPIO11 & O & SPI0 SCLK $\to$ ILI9341 &
    \fixed\ by SoC (SPI0 clock).
    33\,$\Omega$ series resistor on PCB.
    \jref{J1-Pin11} \\

  \rowcolor{pwrcol}
  25 & \na & GND & Ground &
    Additional odd-side ground.
    \jref{J1-Pin15 shared} \\

  \rowcolor{pwrcol}
  27 & GPIO0 & \na & PIN REMOVED --- DO NOT USE &
    HAT ID EEPROM pin. Physical pin removed from header on this prototype. \\

  \rowcolor{oddcol}
  29 & GPIO5 & \na & Spare / SD Card CS &
    Unassigned. May be used as active-low chip select for
    microSD slot on 1651 shield if SD card use is implemented.
    33\,$\Omega$ series resistor on PCB.
    \jref{J1-Pin8} \\

  \rowcolor{rowB}
  31 & GPIO6 & O & TFT DC (Data/Command) &
    Data/Command select line for ILI9341.
    High = data byte; Low = command byte.
    Driven by software before each SPI0 transaction.
    33\,$\Omega$ series resistor on PCB.
    \jref{J1-Pin7} \\

  \rowcolor{oddcol}
  33 & GPIO13 & PWM & Backlight PWM $\to$ 1651 LITE &
    \textbf{Hardware PWM channel~1 (PWM1).}
    Controls display backlight brightness via 1651 LITE pin.
    Solder the LITE/\#3 jumper on the 1651 shield to enable.
    33\,$\Omega$ series resistor on PCB.
    See Section~\ref{sec:pwm} for shared clock constraints with GPIO12.
    \jref{J1-Pin9} \\

  \rowcolor{rowB}
  35 & GPIO19 & I & SPI1 MISO (breakout) &
    SPI1 MISO pin. Not currently connected to any active device.
    Exposed at J4 for future use (e.g.\ a second SPI1 peripheral with MISO).
    Note: GPIO19 is also hardware PWM capable --- not used for PWM here.
    \jref{J4-Pin4} \\

  \rowcolor{pwrcol}
  37 & GPIO26 & \na & Spare &
    Unassigned. Routed to J1 for future use.
    \jref{J1-Pin16} \\

  \rowcolor{pwrcol}
  39 & \na & GND & Ground &
    Odd-side ground. \\

  \midrule
  % ── EVEN SIDE ───────────────────────────────────────────────
  \multicolumn{5}{l}{\cellcolor{evencol}%
    \textbf{Even-side pins --- general I/O: MOSFETs, SPI1 shift register,
    buzzer, 1-Wire, interface headers}} \\
  \midrule

  \rowcolor{pwrcol}
  2 & \na & PWR & 5\,V $\to$ 1651 VIN &
    Powers Adafruit 1651 shield via J1.
    Shield's onboard regulator converts to 3.3\,V for ILI9341 and TSC2007.
    \textbf{Do not use RPi 3.3\,V rail for 1651 power.}
    100\,nF ceramic + 10\,$\mu$F tantalum decoupling at J1 (fitted).
    \jref{J1-Pin1} \\

  \rowcolor{pwrcol}
  4 & \na & PWR & 5\,V $\to$ LED loads &
    Supplies 5\,V to LED/MOSFET drain circuits via J5--J8 Pin~1.
    100\,nF ceramic decoupling near each LED bank (fitted).
    \jref{J5-Pin1, J6-Pin1, J7-Pin1, J8-Pin1} \\

  \rowcolor{pwrcol}
  6 & \na & GND & Ground &
    Even-side primary ground return. \\

  \rowcolor{evencol}
  8 & GPIO14 & \na & UART TX --- DO NOT USE &
    Linux serial console TX. Leave unconnected. \\

  \rowcolor{rowB}
  10 & GPIO15 & \na & UART RX --- DO NOT USE &
    Linux serial console RX. Leave unconnected. \\

  \rowcolor{evencol}
  12 & GPIO18 & \na & Dual-purpose via jumpers &
    GPIO18 routes to two destinations via solder jumpers.
    \textbf{JP2 (closed):} GPIO18 $\to$ J3-Pin4 (interface header).
    \textbf{JP1 (open):} GPIO18 $\to$ J4-Pin6 (interface header).
    Current state: JP2 closed, JP1 open.
    GPIO18 is also SPI1~CE0 --- see Section~\ref{sec:spi1config}.
    \textbf{Warning:} closing both JP1 and JP2 simultaneously
    shorts J3-Pin4 to J4-Pin6. Do not close both jumpers.
    \jref{JP1, JP2, J3-Pin4, J4-Pin6} \\

  \rowcolor{pwrcol}
  14 & \na & GND & Ground &
    Additional even-side ground. \\

  \rowcolor{evencol}
  16 & GPIO23 & \na & Spare / JP2 pin~3 &
    Routed to JP2 pin~3. Not yet dedicated to a function.
    Available for button, indicator, or interface signal.
    \jref{JP2-Pin3, J3-Pin3} \\

  \rowcolor{rowB}
  18 & GPIO24 & O & 74HC595 RCLK (latch) &
    Output. Drives 74HC595 pin~12 (RCLK / storage register clock).
    A rising edge on this pin latches the shift register contents
    to the 74HC595 output pins (QA--QH), updating all 8 MOSFET gates
    simultaneously.
    Pulse this pin \emph{after} each SPI1 byte transmission to update outputs.
    33\,$\Omega$ series resistor on PCB. \\

  \rowcolor{pwrcol}
  20 & \na & GND & Ground &
    Additional even-side ground. \\

  \rowcolor{evencol}
  22 & GPIO25 & O & MOSFET channel~10 (direct GPIO) &
    Drives gate of AO3400A MOSFET for LED channel~10.
    500\,$\Omega$ gate resistor + 10\,k$\Omega$ bleed resistor to GND.
    See Section~\ref{sec:mosfet} for switching caution. \\

  \rowcolor{rowB}
  24 & GPIO8 & O & TFT CS $\to$ ILI9341 (SPI0 CE0) &
    \fixed\ by SoC (SPI0 CE0). Hardware-managed chip select.
    Asserted low automatically by SPI0 hardware during display transactions.
    33\,$\Omega$ series resistor on PCB.
    \jref{J1 ICSP area} \\

  \rowcolor{evencol}
  26 & GPIO7 & O & MOSFET channel~9 (direct GPIO) &
    Drives gate of AO3400A MOSFET for LED channel~9.
    500\,$\Omega$ gate resistor + 10\,k$\Omega$ bleed resistor to GND.
    Note: GPIO7 is SPI0~CE1 --- not a conflict since SPI0~CE1 is unused.
    See Section~\ref{sec:mosfet} for switching caution. \\

  \rowcolor{pwrcol}
  28 & GPIO1 & \na & PIN REMOVED --- DO NOT USE &
    HAT ID EEPROM pin. Physical pin removed on this prototype. \\

  \rowcolor{pwrcol}
  30 & \na & GND & Ground &
    Additional even-side ground. \\

  \rowcolor{rowB}
  32 & GPIO12 & PWM & Buzzer (via MMBT3904 BJT) &
    \textbf{Hardware PWM channel~0-A (PWM0).}
    PWM signal drives base of MMBT3904 NPN BJT.
    BJT collector connects to buzzer ($-$) terminal.
    Buzzer ($+$) connects to 5\,V rail.
    Flyback diode fitted: cathode (line-marked end) to buzzer($+$)/5\,V,
    anode to buzzer($-$)/collector.
    1\,k$\Omega$ base resistor.
    See Section~\ref{sec:pwm} for shared clock constraints with GPIO13. \\

  \rowcolor{pwrcol}
  34 & \na & GND & Ground &
    Additional even-side ground. \\

  \rowcolor{evencol}
  36 & GPIO16 & IO & 1-Wire data bus &
    \textbf{Non-default 1-Wire pin.}
    Requires \texttt{dtoverlay=w1-gpio,gpiopin=16} in
    \texttt{/boot/firmware/config.txt} (BCM number, not physical pin).
    4.7\,k$\Omega$ pull-up to 3.3\,V on PCB (hard requirement of 1-Wire protocol;
    internal RPi pull-up is too weak).
    Connects to J2-Pin2 (1-Wire probe plug) and J5--J8 Pin~8 (LED plugs).
    100\,nF ceramic decoupling at each device end of J5--J8 (fitted).
    See Section~\ref{sec:1wire} for J2 cable note.
    \jref{J2-Pin2, J5-Pin8, J6-Pin8, J7-Pin8, J8-Pin8} \\

  \rowcolor{rowB}
  38 & GPIO20 & O & SPI1 MOSI $\to$ 74HC595 SER &
    \fixed\ for SPI1 function. Drives 74HC595 pin~14 (SER / serial data input).
    Also exposed at J4-Pin5 for potential future SPI1 device.
    33\,$\Omega$ series resistor on PCB.
    \jref{74HC595 pin~14, J4-Pin5} \\

  \rowcolor{pwrcol}
  39 & \na & GND & Ground &
    Odd-side ground (between pins 38 and 40). \\

  \rowcolor{evencol}
  40 & GPIO21 & O & SPI1 SCLK $\to$ 74HC595 SRCLK &
    \fixed\ for SPI1 function. Drives 74HC595 pin~11 (SRCLK / shift register clock).
    Each rising edge shifts one bit from SER into the shift register.
    Also exposed at J4-Pin3 for potential future SPI1 device.
    33\,$\Omega$ series resistor on PCB.
    \jref{74HC595 pin~11, J4-Pin3} \\

\end{longtable}

% =============================================================
\section{74HC595 Shift Register --- As-Built Pin Table}
\label{sec:595}

The 74HC595 provides 8 parallel output lines (QA--QH) from a 3-wire
SPI1 serial interface.  It controls LED channels 1--8 via AO3400A MOSFETs.

\vspace{4pt}
\setlength{\tabcolsep}{4pt}
\renewcommand{\arraystretch}{1.3}

\begin{longtable}{C{1.2cm} C{1.4cm} C{1.6cm} L{3.2cm} L{6.8cm}}
  \toprule
  \rowcolor{hdr}
  \color{white}\textbf{595 Pin} &
  \color{white}\textbf{Name} &
  \color{white}\textbf{Connects to} &
  \color{white}\textbf{Function} &
  \color{white}\textbf{Notes} \\
  \midrule
  \endfirsthead
  \toprule
  \rowcolor{hdr}
  \color{white}\textbf{595 Pin} &
  \color{white}\textbf{Name} &
  \color{white}\textbf{Connects to} &
  \color{white}\textbf{Function} &
  \color{white}\textbf{Notes} \\
  \midrule
  \endhead
  \bottomrule
  \endfoot

  \rowcolor{rowB}
  1 & QB & MOSFET ch~2 & Parallel output B &
    500\,k$\Omega$ to MOSFET gate, 10\,k$\Omega$ bleed to GND.
    LED channel~2. \\
  \rowcolor{oddcol}
  2 & QC & MOSFET ch~3 & Parallel output C & LED channel~3. \\
  \rowcolor{rowB}
  3 & QD & MOSFET ch~4 & Parallel output D & LED channel~4. \\
  \rowcolor{oddcol}
  4 & QE & MOSFET ch~5 & Parallel output E & LED channel~5. \\
  \rowcolor{rowB}
  5 & QF & MOSFET ch~6 & Parallel output F & LED channel~6. \\
  \rowcolor{oddcol}
  6 & QG & MOSFET ch~7 & Parallel output G & LED channel~7. \\
  \rowcolor{rowB}
  7 & QH & MOSFET ch~8 & Parallel output H & LED channel~8. \\
  \rowcolor{pwrcol}
  8 & GND & GND & Ground &
    \textbf{Caution:} 100\,nF decoupling cap fitted at VCC (pin~16),
    but ground path to pin~8 is circuitous.
    Solder a short jumper wire directly from cap ground pad to pin~8
    before testing to improve decoupling effectiveness. \\
  \rowcolor{rowB}
  9 & QH' & Not connected & Serial cascade output &
    Not used. No daisy-chain in this design. \\
  \rowcolor{oddcol}
  10 & !SRCLR & 3.3\,V & Shift register clear &
    Tied high (3.3\,V). Active-low --- holding high keeps shift register enabled.
    To clear all outputs: pulse low in software (requires a GPIO if needed later). \\
  \rowcolor{rowB}
  11 & SRCLK & RPi pin~40 / GPIO21 & Shift register clock &
    SPI1 SCLK. Rising edge shifts one bit from SER into the register.
    8 rising edges load a complete byte. \\
  \rowcolor{oddcol}
  12 & RCLK & RPi pin~18 / GPIO24 & Storage register clock (latch) &
    Manual GPIO output (not a hardware SPI signal).
    Rising edge copies all 8 shift register bits to output pins simultaneously.
    Pulse \emph{after} each SPI1 byte write.
    This separation prevents output glitching during bit-shifting. \\
  \rowcolor{rowB}
  13 & !OE & GND & Output enable &
    Tied to GND (active-low). Outputs are permanently enabled.
    To tri-state outputs (e.g.\ power-down): a GPIO and a trace cut would be needed. \\
  \rowcolor{oddcol}
  14 & SER & RPi pin~38 / GPIO20 & Serial data input &
    SPI1 MOSI. Receives one bit per SRCLK rising edge.
    MSB should be sent first for QA=bit7, QH=bit0 (verify with software). \\
  \rowcolor{pwrcol}
  15 & QA & MOSFET ch~1 & Parallel output A &
    First bit clocked in appears here after latch.
    LED channel~1.
    500\,k$\Omega$ to MOSFET gate, 10\,k$\Omega$ bleed to GND. \\
  \rowcolor{pwrcol}
  16 & VCC & 3.3\,V & Supply voltage &
    3.3\,V from RPi rail. Output high level $\approx$3.0--3.3\,V,
    sufficient for AO3400A gate threshold (1.5\,V typical).
    100\,nF ceramic decoupling fitted (see pin~8 note above). \\

\end{longtable}

% =============================================================
\section{Connector Pin Tables}
\label{sec:connectors}

\subsection{J1 --- 16-pin header, Adafruit 1651 shield interface}

Point-to-point wiring between J1 and the 1651 shield header pins.
Series resistors (33\,$\Omega$) are placed on the PCB between the RPi
header and J1 on all signal lines.

\vspace{4pt}
\begin{longtable}{C{1.4cm} C{2.2cm} C{2.5cm} L{7.0cm}}
  \toprule
  \rowcolor{hdr}
  \color{white}\textbf{J1 Pin} &
  \color{white}\textbf{RPi Pin/BCM} &
  \color{white}\textbf{1651 Shield Pin} &
  \color{white}\textbf{Signal / Notes} \\
  \midrule
  \endfirsthead
  \toprule
  \rowcolor{hdr}
  \color{white}\textbf{J1 Pin} &
  \color{white}\textbf{RPi Pin/BCM} &
  \color{white}\textbf{1651 Shield Pin} &
  \color{white}\textbf{Signal / Notes} \\
  \midrule
  \endhead
  \bottomrule
  \endfoot

  \rowcolor{pwrcol}
  1  & Pin~2 / 5\,V    & Arduino 5V / VIN & 5\,V supply to shield. 100\,nF + 10\,$\mu$F decoupling fitted. \\
  \rowcolor{oddcol}
  2  & Pin~3 / GPIO2   & SDA              & I2C SDA $\to$ TSC2007. \\
  \rowcolor{pwrcol}
  3  & GND             & GND              & Ground. \\
  \rowcolor{rowB}
  4  & Pin~5 / GPIO3   & SCL              & I2C SCL $\to$ TSC2007. \\
  \rowcolor{pwrcol}
  5  & GND             & GND              & Ground. \\
  \rowcolor{oddcol}
  6  & \na             & \na              & Not connected. \\
  \rowcolor{rowB}
  7  & Pin~31 / GPIO6  & Arduino D9       & TFT DC (Data/Command select). \\
  \rowcolor{oddcol}
  8  & Pin~29 / GPIO5  & Arduino D4       & SD Card CS (unassigned; spare). \\
  \rowcolor{rowB}
  9  & Pin~33 / GPIO13 & Arduino D3 / LITE & Backlight PWM. Solder LITE/\#3 jumper on shield. \\
  \rowcolor{oddcol}
  10 & Pin~7 / GPIO4   & Arduino D2 / TSIRQ & Touch interrupt (active-low). Solder TSIRQ/\#2 jumper on shield. \\
  \rowcolor{rowB}
  11 & Pin~23 / GPIO11 & ICSP SCK         & SPI0 SCLK. \\
  \rowcolor{oddcol}
  12 & Pin~21 / GPIO9  & ICSP MISO        & SPI0 MISO $\leftarrow$ ILI9341. \\
  \rowcolor{rowB}
  13 & Pin~19 / GPIO10 & ICSP MOSI        & SPI0 MOSI $\to$ ILI9341. \\
  \rowcolor{oddcol}
  14 & Pin~11 / GPIO17 & \na              & Spare GPIO. \\
  \rowcolor{pwrcol}
  15 & GND             & GND              & Ground. \\
  \rowcolor{rowB}
  16 & Pin~37 / GPIO26 & \na              & Spare GPIO. \\

\end{longtable}

\begin{warnbox}
\textbf{1651 VddIO jumper --- must be set before power-on:}
The shield's VddIO solder jumper selects the logic voltage reference for
the onboard 74LVX245 level shifters.
\textbf{Set to 3V} (solder left pad to centre pad).
The RPi Zero 2\,W is a 3.3\,V logic device; incorrect setting risks damaged GPIO pins.

\smallskip
\textbf{Shield jumpers to solder:}
LITE/\#3 (enables backlight PWM control) and TSIRQ/\#2 (enables touch interrupt).
\end{warnbox}

\subsection{J2 --- 3-pin polarised, 1-Wire temperature probe}

\begin{tabular}{C{1.4cm} C{2.2cm} L{9cm}}
  \toprule
  \rowcolor{hdr}
  \color{white}\textbf{J2 Pin} &
  \color{white}\textbf{Signal} &
  \color{white}\textbf{Notes} \\
  \midrule
  1 & 3.3\,V  & Supply to 1-Wire device. \\
  2 & Data    & 1-Wire data. RPi pin~36 / GPIO16. 4.7\,k$\Omega$ pull-up to 3.3\,V on PCB. \\
  3 & GND     & Ground. \\
  \bottomrule
\end{tabular}

\begin{notebox}
\textbf{J2 cable note:} J2 connects to a 1-Wire temperature probe via a
$\approx$1\,metre cable. There is no decoupling capacitor at J2 on the PCB.
The nearest decoupling caps on the 3.3\,V rail are $\approx$5\,inches away
via J5--J8. The 4.7\,k$\Omega$ pull-up resistor is on the PCB side.
This arrangement is expected to work at standard 1-Wire speed (15.4\,kbps),
as the pull-up source impedance is the dominant factor, not proximity of bypass caps.
If communication errors occur during testing, fit a 100\,nF ceramic cap
between pins 1 and 3 of J2 (3.3\,V to GND) as the first corrective step.
\end{notebox}

\subsection{J3 --- 4-pin polarised, interface control (partial)}

\begin{tabular}{C{1.4cm} C{2.2cm} L{9cm}}
  \toprule
  \rowcolor{hdr}
  \color{white}\textbf{J3 Pin} &
  \color{white}\textbf{Signal} &
  \color{white}\textbf{Notes} \\
  \midrule
  1 & GND     & Ground. \\
  2 & 3.3\,V  & 3.3\,V supply. \\
  3 & GPIO23  & RPi pin~16. Spare / unassigned. Probably button or switch. \\
  4 & GPIO18  & RPi pin~12 via \textbf{JP2 (closed)}. See jumper note below. \\
  \bottomrule
\end{tabular}

\subsection{J4 --- 6-pin polarised, interface control / SPI1 breakout}

\begin{tabular}{C{1.4cm} C{2.2cm} L{9cm}}
  \toprule
  \rowcolor{hdr}
  \color{white}\textbf{J4 Pin} &
  \color{white}\textbf{Signal} &
  \color{white}\textbf{Notes} \\
  \midrule
  1 & 3.3\,V  & 3.3\,V supply. \\
  2 & GND     & Ground. \\
  3 & GPIO21  & RPi pin~40. SPI1 SCLK. Also drives 74HC595 SRCLK. \\
  4 & GPIO19  & RPi pin~35. SPI1 MISO. Not currently connected to a device. \\
  5 & GPIO20  & RPi pin~38. SPI1 MOSI. Also drives 74HC595 SER. \\
  6 & GPIO18  & RPi pin~12 via \textbf{JP1 (open)}. Not currently active. \\
  \bottomrule
\end{tabular}

\begin{notebox}
\textbf{JP1 / JP2 --- solder jumpers for GPIO18 routing:}
GPIO18 is connected to both J3-Pin4 (via JP2) and J4-Pin6 (via JP1).
\textbf{Current state: JP2 closed, JP1 open.}
GPIO18 is therefore active only at J3-Pin4.
\textbf{Do not close both jumpers simultaneously} --- this would short
J3-Pin4 directly to J4-Pin6 through GPIO18.
GPIO18 is also the hardware SPI1~CE0 pin; see Section~\ref{sec:spi1config}.
\end{notebox}

\subsection{J5, J6, J7, J8 --- 8-pin ribbon headers, LED banks}

J5 and J6 carry identical signals (LED channels~1--5).
J7 and J8 carry identical signals (LED channels~6--10).
Duplicate headers allow either plug to connect to the same LED bank.

\begin{tabular}{C{1.4cm} C{2.8cm} L{8.5cm}}
  \toprule
  \rowcolor{hdr}
  \color{white}\textbf{Pin} &
  \color{white}\textbf{Signal} &
  \color{white}\textbf{Notes} \\
  \midrule
  1 & 5\,V (RPi pin~4) & LED anode supply. 100\,nF ceramic decoupling near each bank. \\
  2--6 & LED channels   & MOSFET drain-switched.  See channel assignment table below. \\
  7 & GND              & Ground return for LED loads. \\
  8 & 1-Wire data      & GPIO16 / pin~36. 100\,nF ceramic cap at device end (fitted). \\
  \bottomrule
\end{tabular}

\vspace{6pt}
\textbf{LED channel assignment across headers:}

\begin{tabular}{C{1.4cm} C{2.2cm} C{2.8cm} L{5cm}}
  \toprule
  \rowcolor{hdr}
  \color{white}\textbf{Header pin} &
  \color{white}\textbf{J5, J6 channel} &
  \color{white}\textbf{J7, J8 channel} &
  \color{white}\textbf{Drive source} \\
  \midrule
  \rowcolor{rowB}
  2 & Ch~1  & Ch~6  & 74HC595 QA / QB \\
  \rowcolor{oddcol}
  3 & Ch~2  & Ch~7  & 74HC595 QB / QC \\
  \rowcolor{rowB}
  4 & Ch~3  & Ch~8  & 74HC595 QC / QD \\
  \rowcolor{oddcol}
  5 & Ch~4  & Ch~9  & 74HC595 QD / GPIO7 (direct) \\
  \rowcolor{rowB}
  6 & Ch~5  & Ch~10 & 74HC595 QE / GPIO25 (direct) \\
  \bottomrule
\end{tabular}

\begin{notebox}
Channel assignment across J5--J8 is inferred from the as-built description
(channels~1--5 on J5/J6, channels~6--10 on J7/J8, 8 channels via 595, 2 via direct GPIO).
Verify pin~2--6 assignments against the PCB during testing.
\end{notebox}

% =============================================================
\section{OS Configuration --- \texttt{/boot/firmware/config.txt}}
\label{sec:osconfig}

All settings below are required in \texttt{/boot/firmware/config.txt}
(Raspberry Pi OS Bookworm and later; older releases use \texttt{/boot/config.txt}).
Add or confirm each line before first power-on test.

\subsection{SPI0 --- ILI9341 display}

\begin{alltt}\small
dtparam=spi=on
\end{alltt}

Enables SPI0 (GPIO8--11). Required for ILI9341 display and optional SD card.
CE0 (GPIO8) is managed automatically by the spidev driver.

\subsection{SPI1 --- 74HC595 shift register}
\label{sec:spi1config}

\begin{alltt}\small
dtoverlay=spi1-1cs,cs0_pin=24
\end{alltt}

\begin{warnbox}
\textbf{SPI1 CE0 conflict --- read carefully before configuring.}

The default \texttt{spi1-1cs} overlay assigns \textbf{GPIO18 as CE0}.
GPIO18 in this design is used for JP2/J3 routing --- not for SPI1 CE0.

The 74HC595 does not use a hardware chip select at all.
RCLK (GPIO24) is controlled manually in software to latch outputs.

The recommended approach is:
\begin{enumerate}[noitemsep]
  \item Load the overlay with a remapped CS pin that is not in use:
        \texttt{dtoverlay=spi1-1cs,cs0\_pin=24} redirects CE0 to GPIO24.
        However, GPIO24 is already used for RCLK.
        A cleaner alternative: use \texttt{cs0\_pin=X} where X is a confirmed
        spare GPIO (e.g.\ GPIO23/pin~16 if it remains unassigned).
  \item In Python, open SPI1 with \texttt{no\_cs=True} via spidev,
        disabling hardware CE entirely:
\begin{alltt}\small
import spidev
spi1 = spidev.SpiDev()
spi1.open(1, 0)
spi1.no_cs = True
spi1.max_speed_hz = 1_000_000
\end{alltt}
  \item After each \texttt{spi1.writebytes([byte])}, pulse GPIO24 (RCLK) manually:
\begin{alltt}\small
import RPi.GPIO as GPIO
GPIO.setmode(GPIO.BCM)
GPIO.setup(24, GPIO.OUT)

def shift_out(byte_val):
    spi1.writebytes([byte_val])
    GPIO.output(24, GPIO.HIGH)   # latch rising edge
    GPIO.output(24, GPIO.LOW)    # return low
\end{alltt}
\end{enumerate}
This keeps GPIO18 free and gives precise control of the latch timing.
\end{warnbox}

\subsection{I2C --- TSC2007 touch controller}

\begin{alltt}\small
dtparam=i2c_arm=on
\end{alltt}

Enables I2C bus~1 (GPIO2, GPIO3).
TSC2007 default I2C address: \texttt{0x48}.
Verify with \texttt{sudo i2cdetect -y 1} after wiring.

\subsection{1-Wire --- temperature probe on GPIO16}
\label{sec:1wire}

\begin{alltt}\small
dtoverlay=w1-gpio,gpiopin=16
\end{alltt}

\begin{warnbox}
\textbf{Non-default pin.} The standard 1-Wire pin is GPIO4.
This design uses \textbf{GPIO16 / physical pin~36}.
The overlay parameter \texttt{gpiopin=} takes the \textbf{BCM number (16)},
not the physical pin number.
Using \texttt{gpiopin=7} (the BCM number of physical pin~26, SPI0~CE1) would
configure the wrong pin and leave GPIO16 inactive.

\smallskip
The 4.7\,k$\Omega$ pull-up to 3.3\,V on the PCB is a hard requirement.
The RPi internal pull-up ($\approx$50\,k$\Omega$) is insufficient for 1-Wire timing.

\smallskip
After configuration, probe devices will appear under:
\texttt{/sys/bus/w1/devices/}
Temperature sensors (DS18B20 etc.) register as \texttt{28-xxxxxxxxxxxx}.
Read temperature:
\texttt{cat /sys/bus/w1/devices/28-*/w1\_slave}
\end{warnbox}

\subsection{Hardware PWM --- GPIO12 (buzzer) and GPIO13 (backlight)}
\label{sec:pwm}

Hardware PWM requires no \texttt{config.txt} entry --- it is always available
on GPIO12, GPIO13, GPIO18, GPIO19.
However, using it from Python requires either \texttt{pigpio} or direct
register access. \textbf{Use \texttt{pigpio} --- it handles all clock arithmetic.}

\begin{alltt}\small
# Start the pigpio daemon (add to /etc/rc.local or use systemd):
sudo pigpiod

# Python usage:
import pigpio
pi = pigpio.pi()

# Set backlight first, then buzzer
# (pigpio optimises the shared clock for the last call;
#  setting buzzer last gives it priority for exact frequency)
pi.hardware_PWM(13, 2000, 700000)   # backlight: 2000 Hz, 70% duty
pi.hardware_PWM(12, 2048, 500000)   # buzzer: 2048 Hz, 50% duty

# Adjust buzzer frequency in real time (fine-grained, 1 Hz steps):
for freq in range(1500, 2500, 10):
    pi.hardware_PWM(12, freq, 500000)
    time.sleep(0.05)

# Silence buzzer without disturbing backlight:
pi.hardware_PWM(12, 2048, 0)        # duty=0 silences output

# Adjust backlight brightness (0-100%):
def set_backlight(percent):
    pi.hardware_PWM(13, 2000, percent * 10000)
\end{alltt}

\begin{warnbox}
\textbf{Shared PWM clock constraint --- GPIO12 and GPIO13.}

GPIO12 and GPIO13 share the same hardware PWM clock divider (PWM0/PWM1 on the RP3A0).
Their output frequencies are derived from a single base clock frequency.
This means:
\begin{itemize}[noitemsep]
  \item You can set \emph{independent duty cycles} on GPIO12 and GPIO13.
  \item You \emph{cannot} set \emph{completely independent frequencies} simultaneously ---
        both frequencies must be achievable from the same base clock.
  \item \texttt{pigpio} handles this automatically: it finds the best shared clock
        for both requested frequencies and minimises error.
        Typical frequency accuracy is within 1\% across the audio range.
  \item \textbf{Always set the backlight first, then the buzzer.}
        \texttt{pigpio} optimises the clock for the last call,
        giving the buzzer (the frequency-critical channel) priority.
  \item At 2000\,Hz (backlight) and 2048\,Hz (buzzer), \texttt{pigpio} will find
        a common base clock that satisfies both to within a few Hz.
        The backlight frequency may deviate slightly from 2000\,Hz ---
        this is imperceptible and does not affect LED brightness control.
\end{itemize}

\textbf{Practical buzzer frequency range for CEM-1203(42):}
\begin{itemize}[noitemsep]
  \item Rated resonant frequency: 2048\,Hz (loudest output)
  \item Useful audible range: 800--4000\,Hz
  \item Best efficiency / loudest: 1500--2500\,Hz
  \item \texttt{pigpio} supports any integer Hz value; 1\,Hz steps throughout the range
\end{itemize}
\end{warnbox}

\subsection{Complete \texttt{config.txt} additions summary}

Add the following lines to \texttt{/boot/firmware/config.txt}:

\begin{alltt}\small
# Colorimeter prototype configuration
dtparam=spi=on                        # SPI0: ILI9341 display
dtparam=i2c_arm=on                    # I2C1: TSC2007 touch
dtoverlay=w1-gpio,gpiopin=16          # 1-Wire: GPIO16 (non-default)
dtoverlay=spi1-1cs,cs0_pin=X         # SPI1: 74HC595 (replace X with spare GPIO)
\end{alltt}

Enable and start \texttt{pigpio} daemon for hardware PWM:
\begin{alltt}\small
sudo systemctl enable pigpiod
sudo systemctl start pigpiod
\end{alltt}

% =============================================================
\section{Hardware Cautions for Testing}
\label{sec:cautions}

\subsection{74HC595 gate resistors --- slow MOSFET switching}
\label{sec:mosfet}

\begin{notebox}
The 74HC595 output pins connect to AO3400A MOSFET gates through
\textbf{500\,k$\Omega$} resistors (10\,k$\Omega$ bleed resistors to GND).
The direct-GPIO MOSFETs (channels~9 and~10) use 500\,$\Omega$ gate resistors.

\smallskip
\textbf{Effect of 500\,k$\Omega$ gate resistor:}
The AO3400A input capacitance ($C_{iss}$) is approximately 130\,pF.
Time constant $\tau = 500\,\text{k}\Omega \times 130\,\text{pF} = 65\,\mu$s.
The MOSFET reaches full enhancement (Vgs $\approx$ 3.3\,V) after approximately
$5\tau = 325\,\mu$s.
Turn-off time through the 10\,k$\Omega$ bleed is approximately
$5 \times 10\,\text{k}\Omega \times 130\,\text{pF} = 6.5\,\mu$s (faster).

\smallskip
\textbf{Practical impact:}
For LED switching at 100--5000\,ms intervals, the 325\,$\mu$s turn-on delay
is negligible ($<$0.3\% of the shortest cycle).
The MOSFET will briefly pass through its linear region during turn-on,
dissipating some heat.
At your duty cycles ($<$10 cycles per analysis, $<$1 analysis per 10 minutes),
this is not a thermal concern.

\smallskip
\textbf{Observation to watch for during testing:}
LED onset may appear slightly soft (not an instantaneous step) on an oscilloscope.
If any MOSFETs feel warm to the touch after repeated cycling, note it for revision.
The AO3400A is rated 5.7\,A continuous drain current --- it is not at risk of damage.
\end{notebox}

\subsection{74HC595 decoupling ground path}

The 100\,nF decoupling capacitor at the 74HC595 VCC pin has a circuitous
ground path back to pin~8 (GND).
\textbf{Before testing:} solder a short jumper wire directly from the capacitor's
ground pad to pin~8 of the 74HC595 to provide a low-inductance local ground return.
Without this, the decoupling is partially ineffective at SPI clock frequencies.

\subsection{SPI1 CE0 and GPIO18}

See Section~\ref{sec:spi1config}.
The \texttt{spi1-1cs} overlay claims GPIO18 as CE0 by default.
If the overlay is loaded without the \texttt{cs0\_pin} parameter,
GPIO18 will be driven by the kernel SPI driver, which will conflict
with the JP2 routing to J3-Pin4.
Configure the overlay correctly \emph{before} first power-on with SPI1 active.

\subsection{Buzzer flyback diode orientation}

The flyback diode is fitted with the cathode (line-marked end) at buzzer(+) / 5\,V,
and the anode at buzzer($-$) / MMBT3904 collector.
This is the correct orientation for inductive kickback suppression:
when the BJT switches off, the collapsing magnetic field in the buzzer coil
drives the collector positive, forward-biasing the diode and clamping
the voltage spike.
\textbf{Verify diode orientation before first power-on.}

% =============================================================
\section{Adafruit 1651 Shield Pin Mapping}
\label{sec:1651}

\vspace{4pt}
\renewcommand{\arraystretch}{1.30}

\begin{longtable}{C{1.6cm} C{1.5cm} L{2.8cm} L{6.8cm}}
  \toprule
  \rowcolor{hdr}
  \color{white}\textbf{1651 Pin} \newline \color{white}\footnotesize(Arduino label) &
  \color{white}\textbf{RPi Pin / BCM} &
  \color{white}\textbf{Description} &
  \color{white}\textbf{Notes} \\
  \midrule
  \endfirsthead
  \toprule
  \rowcolor{hdr}
  \color{white}\textbf{1651 Pin} \newline \color{white}\footnotesize(Arduino label) &
  \color{white}\textbf{RPi Pin / BCM} &
  \color{white}\textbf{Description} &
  \color{white}\textbf{Notes} \\
  \midrule
  \endhead
  \bottomrule
  \endfoot

  \rowcolor{pwrcol}
  \textbf{5V / VIN} & Pin~2 / 5\,V &
  Board power in &
  5\,V from RPi. Onboard regulator supplies 3.3\,V to ILI9341 and TSC2007. \\

  \rowcolor{pwrcol}
  \textbf{GND} & GND &
  Ground &
  J1-Pin3, Pin5, Pin15. \\

  \rowcolor{oddcol}
  \textbf{ICSP MOSI} & Pin~19 / GPIO10 &
  SPI0 MOSI $\to$ ILI9341 &
  \fixed. J1-Pin13. 33\,$\Omega$ series resistor. \\

  \rowcolor{rowB}
  \textbf{ICSP MISO} & Pin~21 / GPIO9 &
  SPI0 MISO $\leftarrow$ ILI9341 &
  \fixed. J1-Pin12. 33\,$\Omega$ series resistor. \\

  \rowcolor{oddcol}
  \textbf{ICSP SCK} & Pin~23 / GPIO11 &
  SPI0 clock &
  \fixed. J1-Pin11. 33\,$\Omega$ series resistor. \\

  \rowcolor{evencol}
  \textbf{D10} --- TFT CS & Pin~24 / GPIO8 &
  Display chip select (SPI0 CE0) &
  \fixed\ (SPI0 CE0). Hardware-managed. 33\,$\Omega$ series resistor. \\

  \rowcolor{rowB}
  \textbf{D9} --- TFT DC & Pin~31 / GPIO6 &
  Data/Command select &
  Soft. J1-Pin7. 33\,$\Omega$ series resistor. \\

  \rowcolor{oddcol}
  \textbf{D4} --- SD CS & Pin~29 / GPIO5 &
  microSD chip select &
  Unassigned. J1-Pin8. Wire present; not yet driven in software. \\

  \rowcolor{rowB}
  \textbf{SDA} & Pin~3 / GPIO2 &
  I2C SDA $\to$ TSC2007 &
  \fixed. Shield has onboard 4.7\,k$\Omega$ pull-up. J1-Pin2. \\

  \rowcolor{oddcol}
  \textbf{SCL} & Pin~5 / GPIO3 &
  I2C SCL $\to$ TSC2007 &
  \fixed. Shield has onboard 4.7\,k$\Omega$ pull-up. J1-Pin4. \\

  \rowcolor{rowB}
  \textbf{D2} --- TSIRQ & Pin~7 / GPIO4 &
  Touch interrupt (active-low) &
  Soft. Solder TSIRQ/\#2 jumper on shield. RPi internal pull-up in software. J1-Pin10. \\

  \rowcolor{oddcol}
  \textbf{D3} --- LITE & Pin~33 / GPIO13 &
  Backlight PWM &
  Soft. Solder LITE/\#3 jumper on shield. Hardware PWM1. J1-Pin9. \\

  \rowcolor{warn}
  \textbf{VddIO jumper} &
  \multicolumn{3}{L{11.1cm}}{%
    \textbf{CRITICAL:} Set to \textbf{3V} (solder left pad to centre pad).
    Selects 3.3\,V as logic reference for onboard level shifters.
    Incorrect setting (5V) risks permanent GPIO damage.} \\

\end{longtable}

% =============================================================
\section{Spare and Reserved Pins}
\label{sec:spare}

\begin{tabular}{C{1.8cm} C{1.8cm} L{9.5cm}}
  \toprule
  \rowcolor{hdr}
  \color{white}\textbf{Physical} & \color{white}\textbf{BCM} &
  \color{white}\textbf{Status} \\
  \midrule
  \rowcolor{rowB}
  Pin~11 & GPIO17 & Spare. Routed to J1-Pin14. No function assigned. \\
  \rowcolor{oddcol}
  Pin~13 & GPIO27 & No connect. Available. \\
  \rowcolor{rowB}
  Pin~15 & GPIO22 & No connect. Available. \\
  \rowcolor{oddcol}
  Pin~16 & GPIO23 & Routed to JP2-Pin3 and J3-Pin3. Function TBD (button/switch likely). \\
  \rowcolor{rowB}
  Pin~18 & GPIO24 & 74HC595 RCLK (assigned). \\
  \rowcolor{oddcol}
  Pin~29 & GPIO5  & Spare. Routed to J1-Pin8 (SD CS). Not yet driven in software. \\
  \rowcolor{rowB}
  Pin~35 & GPIO19 & SPI1 MISO breakout at J4-Pin4. No device connected. \\
  \rowcolor{oddcol}
  Pin~37 & GPIO26 & Spare. Routed to J1-Pin16. No function assigned. \\
  \rowcolor{pwrcol}
  Pin~27 & GPIO0  & EEPROM. Physical pin removed. DO NOT USE. \\
  \rowcolor{pwrcol}
  Pin~28 & GPIO1  & EEPROM. Physical pin removed. DO NOT USE. \\
  \rowcolor{pwrcol}
  Pin~8  & GPIO14 & UART TX. Leave unconnected. \\
  \rowcolor{pwrcol}
  Pin~10 & GPIO15 & UART RX. Leave unconnected. \\
  \bottomrule
\end{tabular}

\end{document}
